\section{Qualitative processes and compositional models: bounded proofs}
\label{sec:annex-p16-qualitative-process-and-model-reasoning}

Each proof uses the exact domain stated in the reader. None is a proof of
the complete ENVISION, QPT or QSIM algorithm or of a source unit's closure.

\begin{annexitem}{annex:p16.qr-signs}{thm:p16-qr-signs}{Every local sign case and its realization}
\begin{proof}
Adding zero preserves a real number and hence its sign. Two positive numbers
have a positive sum; two negative numbers have a negative sum. For a positive
and negative number, the pairs \((2,-1),(1,-1),(1,-2)\) give respectively
positive, zero and negative sums. Commuting the two operands gives the
opposite-order witnesses. These cases exhaust the nine ordered pairs of
input signs and prove both containment and realization for the displayed table.
For multiplication, a zero factor gives zero, like nonzero signs give a
positive product and unlike signs a negative product. Choosing factors one
or minus one realizes each nonzero case. These are single-operation witnesses;
they do not promise one set of magnitudes simultaneously realizing a network.
\end{proof}
\end{annexitem}

\begin{annexitem}{annex:p16.qr-balance}{thm:p16-qr-balance}{The direct-influence adder and activity distinction}
\begin{proof}
The stipulated complete direct-influence set gives exactly \(E'=u-v\).
Subtracting the two positive real rates yields the three signs according to
their order. The choices \((u,v)=(2,1),(1,1),(1,2)\) realize all cases
with both processes present. In the middle case the derivative is zero,
but the loss activation condition is separately stipulated true and both
rates remain positive. Neither cancellation nor the derivative sign changes
that condition. If these equal rates persist over an interval, the mean
value theorem makes the stock constant there; equality at a single instant
alone makes no such interval claim.
\end{proof}
\end{annexitem}

\begin{annexitem}{annex:p16.qr-monotone}{thm:p16-qr-monotone}{Interior chain rule and range-endpoint cases}
\begin{proof}
At an interior time with \(g(t)\) in the interior of its value range,
the chain rule gives \(f'(t)=H'(g(t))g'(t)\). A positive multiplier
preserves each of the three signs; a negative one reverses the two nonzero
signs and preserves zero. If \(g(t)\) is an endpoint of its range, it
is a local extremum at this interior time, so Fermat's theorem gives
\(g'(t)=0\). Monotonicity makes \(f(t)=H(g(t))\) a local extremum
too, and its declared differentiability gives \(f'(t)=0\). This argument
does not require a nonzero derivative of \(H\) at a range endpoint.
If \(g\) is constant, then so is \(f\), and both signs are zero.
These cases exhaust finite real-valued ranges of a continuous function on
the compact interval. For \(H(x)=x^3,g(t)=t\), direct differentiation
gives \(f'(0)=0\) while \(g'(0)=1\). If an additional input varies,
for example \(f=g-h\), the additional derivative \(-h'\) enters;
the one-input chain-rule premises no longer hold.
\end{proof}
\end{annexitem}

\begin{annexitem}{annex:p16.qr-valve}{thm:p16-qr-valve}{One confluence, not a global device solution}
\begin{proof}
The confluence's admissible tuples are abstractions of
\(dP+dA-dQ=0\). With \(dQ=0\) and \(dA>0\), this becomes
\(dP=-dA<0\); the result is forced for every realizing pair of magnitudes.
With \(dP>0,dA<0\), the equation gives \(dQ=dP+dA\).
The witnesses in the sign-table proof realize all three possible flow-change
signs. Thus the second case is ambiguous. This argument uses no other
component or mode and says nothing about realizability of an entire device.
\end{proof}
\end{annexitem}
\begin{annexitem}{annex:p16.qr-containers}{thm:p16-qr-containers}{Every initial direction and the reached-boundary argument}
\begin{proof}
Positive \(q_F,q_H\) give \(V_F',V_H'>0,V_G'<0\).
The chain rule gives \(h_i'=\beta_iV_i'\), so the three head signs
are the same. Consequently
\[
 (h_G-h_F)'=-\beta_G(q_F+q_H)-\beta_Fq_F<0,
 \quad (h_G-h_H)'=-\beta_G(q_F+q_H)-\beta_Hq_H<0.
\]
Adding the stock derivatives gives zero; the mean value theorem makes
total amount constant on the interval. These signs contain no specified
initial head gaps or rate magnitudes and thus do not determine an ordering
of the two contact times.

At the stated reached boundary, \(q_F=0,q_H>0\). Hence
\((h_G-h_F)'=-\beta_Gq_H<0\) there. By right-continuity of this
derivative it remains negative on some sufficiently short right interval.
Integrating from the zero gap gives a negative gap at each sufficiently
near later time. The retained pressure relation and open path then enable
reverse flow. Exchanging names proves the symmetric result. If both gaps
reach zero simultaneously, this particular positive-other-flow premise
fails and the reversal conclusion is not asserted. External actions,
different elevations or pressure contributions are outside this proof.
\end{proof}
\end{annexitem}

\begin{annexitem}{annex:p16.qr-contact}{thm:p16-qr-contact}{Equality crossing, a nonhitting limit and a sufficient bound}
\begin{proof}
The gap \(h=x-l\) is continuous. Negative and positive values at two
times force a zero between them by the intermediate value theorem.
For \(x=l-e^{-t}\), the gap is always negative and its derivative
\(e^{-t}\) always positive at finite times, so no finite contact follows.
For the sufficient bound put \(d=l(t_0)-x(t_0)>0\). If no contact
occurred by \(t_0+d/\epsilon\), the stipulated derivative bound and
the fundamental theorem of calculus would give
\(h(t_0+d/\epsilon)\ge -d+\epsilon(d/\epsilon)=0\), contradicting
continued negativity (or giving contact at that endpoint). Validity of
the derivative bound and mode through that interval is a premise.
Finally \(x(t)=t^2\) has derivative zero at zero but is not constant
on any right neighborhood. Thus point stationarity is not an interval result.
\end{proof}
\end{annexitem}

\begin{annexitem}{annex:p16.qr-filter}{thm:p16-qr-filter}{Termination, preservation and a locally consistent obstruction}
\begin{proof}
Each nontrivial step removes a member of one of finitely many finite tuple
sets, so at most their total initial cardinality deletions occur. A completed
fair scan with no deletion is a fixed point; the finite scan terminates.
Let \(s\) be a global assignment satisfying all initial constraints.
Its projection onto any constraint scope has matching support in each
neighbor: the neighbor's projection of the same \(s\). Inductively all
these projections remain in their tuple sets, so none can be deleted.
This proves preservation of every such assignment.

For the triangle, each inequality relation on \(\{0,1\}\) contains
the tuples \((0,1),(1,0)\). For either tuple, each shared variable's
value appears in a supporting neighbor tuple; no deletion occurs. But
\(X\ne Y,Y\ne Z\) on this binary domain implies \(X=Z\), contrary
to \(Z\ne X\). Thus local supports do not ensure a global assignment.
The proof concerns finite label constraints; it does not establish the
validity of every QSIM generator, filter or physical abstraction.
\end{proof}
\end{annexitem}

\begin{annexitem}{annex:p16.qr-energy}{thm:p16-qr-energy}{Conservation and the return comparison}
\begin{proof}
Continuity of \(H\) makes its integral a differentiable potential with
derivative \(H\). The chain rule and the two stated equations give
\[
 {\cal E}'=vv'+H(x)x'=v(-H(x))+H(x)v=0.
\]
The mean value theorem makes \({\cal E}\) constant. At both zero-position
times its integral term is zero, so \(v(t_1)^2=v(t_0)^2\).
Positivity at both returns selects the same square root and gives equality.
No theorem of global periodicity, existence of a return, damping or
termination of a qualitative simulator was needed or inferred.
\end{proof}
\end{annexitem}
